\section{Estimated preprocessing and validation selection}
\label{sec:practical}

We next examine whether the effects in the controlled comparisons remain visible
when we sample training data, estimate preprocessing, and select configurations
by validation. We use a finite grid of candidates, so this experiment does not
extend the population theorems to estimated scales. The full protocol, results,
and reproducible artifacts are recorded in \texttt{notes/practical-selection.md}
and \texttt{notes/practical-selection-results.md}.

For the synthetic experiment, we use eight feasible binary laws with 128
independent replications per law, 256 training observations, and 128 independent
validation observations. We include independent and nested signals at $q = 0.02$
and $q = 0.002$, signals with opposite signs under nested and disjoint
dependence, and rare-only and common-only signals. We compare known
preprocessing with estimated scales, training centering with a fixed or fitted
intercept, and uncentered predictors with a fitted intercept. We estimate
variances with denominator $n$ and give constant training columns divisor one
and coefficient zero. Absent features therefore cannot acquire a fitted effect
in the estimated preprocessing arms. Fitted intercepts start at the training
response mean, while the population-oracle arms start at the true centered
intercept of zero. SGD samples with replacement from a fixed noisy training
set, so revisiting an observation also reuses its noise.

The divisors are $1$, $\sqrt{v_j}$, and $\sqrt{v_j + 0.01}$. For each scaling
and preprocessing arm, we use five rates $\eta = a/\max(L,10^{-12})$,
$a \in \{0.001,0.01,0.05,0.2,0.8\}$, and two batch sizes. Here $L$ is the
largest empirical curvature, including the unscaled intercept coordinate when
fitted. This calibration does not guarantee stochastic stability. We use batches
$B \in \{1,16\}$ and observation caps $M \in \{128,512,2048\}$. We retain
initialization and every checkpoint, and compare fixed $a = 0.05$, $B = 16$ with
validation selection of the rate and batch size, or of the rate, batch size,
and stopping count. We select by either aggregate or rare-presence validation
half squared loss. When the rare validation group is empty, its selection rule
falls back to aggregate loss. We use no test outcomes in selection and do not
refit selected models on validation data. Each combination of scaling and
preprocessing uses ten trajectories, costing $10M$ processed observations at
cap $M$, with checkpoints sharing paths.

Differences in how quickly the coefficients are learned remain visible with
estimated preprocessing. For the independent $(1,1)$ signal at $q = 0.02$ and
noise variance three, the fixed configuration gives an SD-minus-none rare excess
half squared loss of $-0.1653$ (paired MCSE $0.0295$) at $M = 128$ and
$+0.0901$ ($0.0697$) at $M = 2048$ with training centering and a fitted
intercept. The early advantage is clear, while the later disadvantage is within
two MCSE of zero. At the latter budget, the rare-only signal favors SD by
$0.2320$ ($0.0233$), whereas the common-only signal favors no scaling by
$0.1272$ ($0.0159$). These observations concern the declared configurations with
sampled data, rather than the exact calibrations of the population comparisons.

Missing groups create a further problem. At $q = 0.002$, the two anchor laws
have no rare validation observations in 106 and 105 of 128 replications.
Validation can then select a finite but explosive rare coefficient because no
rare cases are available to reveal its error. One SD fit with centered
predictors has a rare coefficient of about $-6.719 \times 10^{16}$ and
validation half squared loss $2.086$, yet its exact rare loss over the binary
cells is about $2.257 \times 10^{33}$. We retain this path in the results. In
a separate investigation, we replay every eligible trajectory for the worst
selected replication in each of 24 combinations of law and scaling, using
160-digit decimal arithmetic. All 24 selections remain unchanged, so the
detected floating-point cancellation does not explain the failure. Means
dominated by these extreme finite outcomes do not precisely estimate population
risk. We also retain all 151 failed synthetic endpoints, flagged in the archive.

\paragraph{A sparse categorical task.}

We fit a squared-loss linear predictor to the recorded poisonous label in the
UCI Mushroom dataset \citep{uciMushroom}, using all 22 categorical predictors.
The declared subgroup consists of mushrooms with green spore-print color, which
accounts for 72 of 8124 records. The records are hypothetical descriptions
derived from a field guide, so evaluation on random partitions does not
establish generalization to new species or field data. There are no duplicate
predictor rows. We construct the vocabulary and estimate centering and scales
from training data alone, treat missingness as a category, and map unknown
categories to zero blocks. We store the one-hot design sparsely and implement
centering algebraically. Because the full one-hot design is rank deficient, its
coefficients depend on the optimizer and cannot be interpreted as uniquely
identified effects.

We use 32 random partitions for each pair of training and validation sizes,
$(256,128)$ and $(2048,1024)$, holding out the remaining records for testing.
All candidates share partitions and replacement index streams. We use batches
$B \in \{16,128\}$ and caps $M \in \{512,2048,8192\}$, and do not clip
predictions. Table~\ref{tab:practical} reports results with centered predictors
and a fitted intercept at the largest cap, jointly selecting the rate, batch
size, and stopping count. We separately retain the other preprocessing arms,
shorter caps, fixed stops, all paths, and the search across scaling rules, which
requires three times as much work.

\begin{table}[t]
  \centering
  \small
  \begin{tabular}{llrr}
    \toprule
    Validation objective & Scaling &    Aggregate loss &         Rare loss \\
    \midrule
    Aggregate            & None    & .001754 (.000060) & .012482 (.001386) \\
    Aggregate            & SD      & .004617 (.002125) & .005912 (.001054) \\
    Aggregate            & Damped  & .001175 (.000034) & .004151 (.000465) \\
    Rare                 & None    & .001878 (.000130) & .012465 (.001382) \\
    Rare                 & SD      & .123905 (.075527) & .003158 (.000451) \\
    Rare                 & Damped  & .014966 (.009145) & .003984 (.000495) \\
    \bottomrule
  \end{tabular}
  \caption{
    Mushroom held-out half squared loss (MCSE), with $n = 2048$, 1024 validation
    records, and budget cap $M = 8192$. Each scaling receives the same ten
    trajectories. MCSE describes variation across partitions and optimization
    runs conditional on this dataset. It does not give a population confidence
    interval.
  }
  \label{tab:practical}
\end{table}

Under aggregate selection, the SD-minus-none rare-loss difference is $-0.00657$
(paired MCSE $0.00153$), and the aggregate difference is $+0.002863$
($0.00213$). Selecting SD by rare rather than aggregate validation loss reduces
rare loss by $0.002754$ ($0.000967$), while increasing aggregate loss by
$0.11929$ ($0.0756$). The green coefficient on the original scale averages
$0.4822$, $0.7462$, and $0.7678$ for none, SD, and damped scaling under
aggregate selection. Other categorical coefficients partly offset its larger
contribution. These results show that relative scaling affects prediction under
a finite budget, with the ordering depending on the objective and tuning.

At the smaller sample size, five training partitions and thirteen validation
partitions contain no green observations. At the larger size, aggregate
validation selects the largest rate factor for unscaled training in every
partition, leaving the optimum unresolved. These limits in coverage and search
range preclude a general normalization recommendation. Independent checks
reconstruct every archived validation and test loss, selection, and paired
uncertainty estimate, and replay representative complete paths in transformed
coordinates. We preserve the earlier controlled studies and their reported
failures.
